%% ---------------------------------------------------------------------------
%% Offline Accuracy Is Not Optimization Validity
%% Focused preprint. Overleaf: upload this folder, set main.tex as the root
%% document, compiler pdfLaTeX. Figures build from committed artifacts with
%%     python make_figs.py
%% and are optional -- the document compiles without them (see \iffigures).
%% ---------------------------------------------------------------------------
\documentclass[journal]{IEEEtran}

\usepackage{amsmath,amssymb}
\usepackage{booktabs}
\usepackage{graphicx}
\usepackage{url}
\usepackage[hidelinks]{hyperref}

% Flip to \figurestrue once make_figs.py has produced figures/*.pdf
\newif\iffigures
\figuresfalse

\newcommand{\fmax}{$F_{\mathit{max}}$}
\newcommand{\TODO}[1]{\textcolor{red}{[TODO: #1]}}
\usepackage{xcolor}

\begin{document}

\title{Offline Accuracy Is Not Optimization Validity:\\
Measuring Learned-Reward Collapse in RTL Generation\\
Against Physical Ground Truth}

\author{Mohamed Amine Hallam\thanks{Harbin Institute of Technology (Shenzhen).
Correspondence: hallamamine05@gmail.com}}

\maketitle

\begin{abstract}
Reinforcement learning for register-transfer level (RTL) generation increasingly
optimizes a \emph{learned} model of physical quality, because invoking a
synthesis tool inside the training loop is prohibitively slow. Such a reward is
validated the way any regressor is: held-out accuracy. We show that this
validation is not sufficient, and we measure the gap against real
electronic design automation (EDA) results rather than a synthetic proxy for
them. On a held-out set of thirty digital signal processing accelerators frozen
before training, a learned maximum-frequency predictor is highly accurate on the
behaviour policy it was fitted near---mean error $+4.7$\,MHz, Spearman
$0.958$---and badly invalid on the policy that optimization produces---mean
error $+235.9$\,MHz, with $25$ of $30$ designs pinned at the predictor's own
output ceiling. Tracking the predictor along the optimization path shows the
failure is abrupt rather than gradual: predictions are stable for roughly two
hundred gradient updates and then saturate, and functional correctness degrades
at exactly that point. A causal control isolates the cause. Removing the
physical term while retaining the correctness gate and the trust-region penalty
raises correctness to its highest observed value, $98.3\%$, while reducing
measured frequency $38\%$ below the supervised starting policy: the two
objectives trade off directly. Finally, we show the failure is repairable at
small cost. Adding $43$ real post-implementation measurements drawn from the
region the optimizer reaches reduces predictor error from $+235.9$ to
$+30.7$\,MHz and eliminates saturation on $29$ of $30$ affected cases;
re-optimizing against the repaired predictor improves measured frequency from
$176.7$ to $192.4$\,MHz and correctness from $86.9\%$ to $95.1\%$ using $30\%$
fewer gradient updates. Notably, the repaired predictor scores \emph{worse} by
conventional offline cross-validation than the one it replaces, so a
practitioner selecting on that metric would have kept the broken reward.
\end{abstract}

\begin{IEEEkeywords}
Reinforcement learning, reward modelling, RTL generation, high-level synthesis,
FPGA, design automation, reward overoptimization.
\end{IEEEkeywords}

% ===========================================================================
\section{Introduction}
% ===========================================================================

\IEEEPARstart{L}{arge} language models are now routinely fine-tuned with
reinforcement learning to produce hardware description code that is not merely
correct but physically fast~\cite{chipseek2025,ppartl2025}. The obstacle is
cost. A single synthesis and implementation run takes minutes; a policy-gradient
step needs a reward for every sample in a group, and a training run needs
hundreds of steps. Invoking the EDA tool in the loop is therefore impractical at
the scale at which these methods operate, and the standard resolution is to
train a surrogate: a learned model mapping a design to its predicted maximum
operating frequency, queried in milliseconds.

The surrogate is a regressor, and it is validated as one---held-out
accuracy, rank correlation, leave-one-design-out cross-validation. If those
numbers are good, the reward is taken to be sound. This paper argues that the
inference is invalid, and supplies the measurement that shows why.

The concern is not new in the abstract. Optimizing against an imperfect proxy is
known to degrade true objective performance once optimization pressure is
sufficient, a phenomenon characterized for language-model reward models by
Gao \emph{et al.}~\cite{gao2023scaling}. Their study, by necessity, used a
synthetic gold-standard reward model in place of ground truth, because
collecting enough genuine human preference data was infeasible.

Hardware generation removes that limitation. The ground truth is a synthesis
tool and a physical device. This lets us ask a question that has not previously
been answerable with real measurements: \emph{when a learned physical reward is
optimized, where and how does it stop being true?}

Our contributions are:

\begin{enumerate}
\item \textbf{A measurement of proxy collapse against real EDA.} On a frozen
held-out set, the same predictor has mean error $+4.7$\,MHz on the supervised
policy and $+235.9$\,MHz on the reinforcement-learned policy
(Section~\ref{sec:divergence}). No held-out fit detects this, because the
distribution on which it fails does not exist until optimization creates it.

\item \textbf{The dynamics of the failure.} Evaluating saved checkpoints along
the optimization path shows predictions stable for approximately two hundred
gradient updates followed by abrupt saturation, with functional correctness
degrading simultaneously (Section~\ref{sec:trajectory}). The correctness
regression is not a separate defect; it is the same event.

\item \textbf{A causal control.} Removing the physical reward term while keeping
the correctness gate and the trust-region penalty produces the highest
correctness we observe ($98.3\%$) and frequency $38\%$ \emph{below} the
supervised starting point (Section~\ref{sec:ablation}). The learned physical
signal is doing the work, and the two objectives genuinely conflict.

\item \textbf{A cheap repair, and evidence that offline metrics misrank it.}
Forty-three real measurements from the optimized region reduce error to
$+30.7$\,MHz and improve every downstream outcome, while scoring worse by
leave-one-design-out cross-validation than the predictor it replaces
(Section~\ref{sec:repair}).
\end{enumerate}

\iffigures
\begin{figure*}[t]
\centering
\includegraphics[width=\textwidth]{figures/fig1_trajectory.pdf}
\caption{Predicted and measured \fmax{} along the optimization path, as a
function of cumulative optimizer updates. Left: the original predictor. Its
predictions track measurement for roughly two hundred updates, then saturate at
the output ceiling while the measured value does not follow. Right: the
re-anchored predictor, on which the two curves remain together. Correctness is
overlaid; it degrades only where saturation occurs.}
\label{fig:trajectory}
\end{figure*}
\fi

% ===========================================================================
\section{Setup}
\label{sec:setup}
% ===========================================================================

\subsection{Task and models}

We generate single-input streaming DSP accelerators sharing one interface (a
clock, an active-low reset, an $8$-bit input sample, a $16$-bit registered
output) across five families spanning three structurally distinct circuit
classes: multiply--accumulate datapaths (\texttt{fir}, \texttt{firr}),
polynomial evaluation (\texttt{poly}), recursive filters with genuine feedback
(\texttt{iir}), and multiplier-free comparator networks (\texttt{med}). The
common interface is what makes a uniform executable oracle and a uniform
measurement harness possible.

The base model is RTLCoder-6.7B~\cite{liu2024rtlcoder}, adapted with
LoRA~\cite{hu2022lora} ($r=16$). Supervised fine-tuning provides the starting
policy; policy optimization uses group-relative advantage
estimation~\cite{shao2024deepseekmath} with group size $G=8$, a generation
budget of $1536$ tokens, and a Kullback--Leibler penalty against a
\emph{frozen copy of the supervised policy} rather than the base model.

\subsection{Reward}

The reward gates a learned frequency prediction on functional correctness:
\begin{equation}
r(m) =
\begin{cases}
\mathrm{clip}\!\left(\hat{F}(m),\, 5,\, 500\right) & \text{if } m \text{ is
correct},\\[2pt]
0 & \text{otherwise},
\end{cases}
\label{eq:reward}
\end{equation}
where $\hat{F}$ is a multilayer perceptron over twelve lexical features of the
generated source, trained to predict $\log$ \fmax{} from real
post-implementation results. The clip bounds are applied in log space before
exponentiation; they exist because an earlier unbounded version overflowed to
infinity on one design family. \textbf{The upper bound of this clip becomes
central to what follows.}

Correctness is decided \emph{only} by an executable oracle that compares the
candidate's simulated output stream against a reference implementation under
Icarus Verilog, allowing for pipeline latency. No similarity heuristic
participates.

\subsection{Held-out protocol}

The held-out split was fixed before any model reported here was trained and is
enforced by a single predicate imported by the corpus generator, the
optimization design list, and the predictor's data loader; each refuses to
proceed if a held-out design appears. It distinguishes \emph{interpolation}
(parameters inside the trained range, withheld) from \emph{extrapolation}
(parameters outside it, constructed at evaluation time), giving $30$ held-out
designs---$19$ interpolation, $11$ extrapolation---against $145$ training
designs.

Evaluation samples $n=48$ per design and decides correctness with two stimulus
seeds disjoint from the training seed, each four times longer.

\subsection{Measurement and its limits}

All frequencies come from Vivado 2023.1 targeting a Xilinx Zynq-7020
(\texttt{xc7z020clg400-1}) with a $5$\,ns target period, computed from the
post-implementation worst negative slack as $1000/(T - \mathrm{WNS})$. We state
plainly that this is a \emph{WNS-derived estimate from a single implementation
run}, not a timing-closure result: we do not iterate the constraint to the point
of closure. It is the standard basis for relative comparison in this
literature~\cite{ustun2020accurate}, and all comparisons here are relative,
between policies on identical designs with an identical flow.

The reported per-design frequency weights each distinct implementation by the
number of samples that collapsed onto it. We report two variants: the mean over
\emph{correct} candidates, and a \emph{penalised} mean scoring incorrect samples
$0$\,MHz. The second is the equal-sample-cost number and we lead with it,
because the first cannot distinguish a policy that is fast from one that is fast
only on the rare occasions it is right.

% ===========================================================================
\section{The Predictor Is Accurate Where It Was Fitted and Invalid Where
Optimization Goes}
\label{sec:divergence}
% ===========================================================================

Table~\ref{tab:divergence} pairs the predictor's output with the measured value
for every (policy, design) cell of the held-out set.

\begin{table}[t]
\caption{Predictor error against real post-implementation measurement, by the
policy that generated the candidates. Identical predictor throughout.}
\label{tab:divergence}
\centering
\small
\begin{tabular}{lrrrrc}
\toprule
policy & pred. & measured & bias & MAE & saturated \\
\midrule
base            & 329.2 &  45.6 & $+283.6$ & 283.6 & 10/16 \\
supervised      &  88.3 &  83.6 & $+4.7$   &  20.5 & 1/30 \\
best-of-8       & 154.5 & 155.5 & $-1.0$   &  35.1 & 0/30 \\
\textbf{optimized} & \textbf{433.6} & \textbf{197.7} & $\mathbf{+235.9}$ &
\textbf{240.1} & \textbf{25/30} \\
\bottomrule
\end{tabular}
\\[2pt]
\footnotesize All values in MHz. ``Saturated'' counts designs whose maximum
prediction sits on the $500$\,MHz clip bound of Eq.~\eqref{eq:reward}.
\end{table}

Three observations follow.

\textbf{The predictor is genuinely good where it was fitted.} On the supervised
policy its bias is $+4.7$\,MHz against measurements spanning $15$--$347$\,MHz,
with Spearman correlation $0.958$. By any conventional standard this is a
successful regressor, and it is the standard by which such rewards are accepted.

\textbf{It is not good where optimization takes the policy.} On the optimized
policy the bias is $+235.9$\,MHz and $25$ of $30$ designs are pinned at the clip
bound. Once predictions pile onto a ceiling, the reward can no longer
distinguish a fast design from a faster one, and the gradient stops carrying
frequency information.

\textbf{The improvement is real; its magnitude is not.} Measured frequency does
rise from $83.6$ to $197.7$\,MHz. The predictor claims a rise from $88.3$ to
$433.6$. The excess,
\begin{equation}
\underbrace{(433.6-88.3)}_{\text{predicted gain}} -
\underbrace{(197.7-83.6)}_{\text{measured gain}} = +231.2~\text{MHz},
\end{equation}
is inflation attributable purely to the proxy. Reward overoptimization here does
not destroy the objective, it overstates progress roughly threefold---which is
harder to notice than outright failure.

\subsection{The error is also family-dependent}

The failure is not solely a function of optimization pressure. On the
\texttt{iir} family the predictor is badly wrong \emph{even on the supervised
policy} ($+154.9$ and $+182.2$\,MHz on two designs, $-86.1$ on a third). These
are precisely the families where optimization delivered least benefit. Where the
proxy is valid the method works; where it is not, it does not. This is a
falsifiable prediction the repair experiment later confirms.

% ===========================================================================
\section{Dynamics: The Failure Is Abrupt, and Correctness Fails With It}
\label{sec:trajectory}
% ===========================================================================

Endpoint measurements cannot distinguish gradual drift from sudden collapse. We
therefore evaluated saved checkpoints along the optimization path with identical
prompts, oracle and measurement flow (Table~\ref{tab:trajectory}).

\begin{table}[t]
\caption{Predictor output and functional correctness along the optimization
path.}
\label{tab:trajectory}
\centering
\small
\begin{tabular}{lrrr}
\toprule
checkpoint & updates & correctness & predicted \fmax{} \\
\midrule
supervised & 0   & 94.3\% &  82.9 \\
c1         & 77  & 95.8\% & 218.1 \\
c2         & 138 & 95.3\% & 213.7 \\
c3         & 205 & 95.3\% & 216.6 \\
\textbf{c4} & \textbf{276} & \textbf{88.0\%} & \textbf{497.9} \\
\bottomrule
\end{tabular}
\end{table}

The predictor rises once, early, to roughly $215$\,MHz and then sits there for
some two hundred gradient updates. Between the third and fourth checkpoint it
jumps to $497.9$---the clip bound of Eq.~\eqref{eq:reward} to within rounding.

Correctness collapses over the same interval, from $95.3\%$ to $88.0\%$, having
been stable or improving until then.

We draw the inference explicitly: \textbf{the correctness regression and the
predictor saturation are one event}, not two. Practitioners who observe a late
correctness regression in correctness-gated optimization may be observing proxy
saturation rather than an independent problem with the gate, and the diagnostic
is cheap---saturation is visible in the training log without any EDA invocation.

% ===========================================================================
\section{Causal Control: The Physical Term Is Load-Bearing
\label{sec:ablation}}
% ===========================================================================

That measured frequency rises during optimization does not establish that the
learned physical reward caused it. Correctness pressure and the trust-region
penalty are also present. We therefore re-ran optimization from the same
supervised checkpoint with Eq.~\eqref{eq:reward} replaced by a constant reward
of $1$ for correct candidates, leaving the correctness gate and the penalty
untouched (Table~\ref{tab:ablation}).

\begin{table}[t]
\caption{Effect of removing the learned physical term. Equal-sample measured
\fmax{} (incorrect samples scored $0$) and correctness.}
\label{tab:ablation}
\centering
\small
\begin{tabular}{lrrrr}
\toprule
policy & interp. & extrap. & all & correctness \\
\midrule
supervised          &  79.9 &  61.2 &  73.0 & 88.7\% \\
correctness-only    &  52.8 &  31.1 &  45.3 & \textbf{98.3\%} \\
physical reward     & \textbf{198.7} & \textbf{138.6} & \textbf{176.7} & 86.9\% \\
\bottomrule
\end{tabular}
\end{table}

Removing the physical term does not merely forgo the gain. Correctness reaches
$98.3\%$, the highest of any policy we trained, while measured frequency falls
$38\%$ below the supervised starting point. The candidate-level view is
sharper than the aggregate: on one design the supervised policy emits
implementations measuring $34$, $43$, $224$, $224$ and $43$\,MHz, while the
correctness-only policy emits $43$, $43$, $43$. The fast implementation is not
made less likely; it is eliminated.

The explanation is straightforward once stated. The implementation most likely
to be correct is the naive direct form, and a purely correctness-driven
objective walks toward it. Speed and correctness are in genuine tension, and the
physical reward is what prevents collapse onto the slowest correct design.

\subsection{A binary reward extinguishes its own gradient}

The control also exposes a practical obstacle to the obvious alternative. Under
group-relative advantage, a group whose rewards are all equal contributes
nothing. With the supervised policy already $\sim\!\!90\%$ correct, most groups
under a binary reward return unanimous: the correctness-only run produced $159$
optimizer updates from $1997$ attempted steps ($92\%$ degenerate), against
$284$ from $408$ ($30\%$) for the physical reward, and drifted $36\times$ less
per update. A binary correctness reward is not merely weaker at this competence
level---it removes the signal it depends on.

% ===========================================================================
\section{Repair, and Why Offline Selection Would Have Missed It}
\label{sec:repair}
% ===========================================================================

If the predictor fails because optimization moves the policy off its training
distribution, then labelling that region should repair it. We synthesised $43$
candidates drawn from the optimized policy on \emph{training-split} designs
only---held-out labels would violate the protocol---spanning $22.0$ to
$338.4$\,MHz, and refitted.

\begin{table}[t]
\caption{Predictor error before and after re-anchoring, scored on identical
stored candidates.}
\label{tab:repair}
\centering
\small
\begin{tabular}{lrrrr}
\toprule
 & \multicolumn{2}{c}{original} & \multicolumn{2}{c}{re-anchored} \\
\cmidrule(lr){2-3}\cmidrule(lr){4-5}
policy & bias & sat. & bias & sat. \\
\midrule
base         & $+283.6$ & 10/16 & $+124.0$ & 2/16 \\
supervised   & $+4.7$   &  1/30 & $+6.8$   & 2/30 \\
optimized    & $+235.9$ & 25/30 & $\mathbf{+30.7}$ & \textbf{2/30} \\
\midrule
\multicolumn{5}{l}{\emph{offline} leave-one-design-out $\rho$:
\;0.965 $\rightarrow$ \textbf{0.841}} \\
\bottomrule
\end{tabular}
\end{table}

Bias on the optimized distribution falls from $+235.9$ to $+30.7$\,MHz and
saturation from $25$ of $30$ designs to $2$. Aggregate saturation across all
cells falls from $36$ of $106$ to $7$.

\textbf{And the repaired predictor is worse by the conventional offline
metric}---leave-one-design-out rank correlation drops from $0.965$ to $0.841$. A
practitioner selecting a reward model by cross-validation, which is the normal
practice, would have kept the broken one. This is the paper's thesis applied to
its own model-selection step.\footnote{The two cross-validation figures are
computed on different pair sets ($203$ and $272$), so this is an observation
rather than a controlled comparison; we report it as such.}

\subsection{Re-optimizing against the repaired predictor}

\begin{table}[t]
\caption{Optimization against the original versus the re-anchored predictor.
Equal-sample measured \fmax{}.}
\label{tab:v9}
\centering
\small
\begin{tabular}{lrr}
\toprule
 & original & re-anchored \\
\midrule
measured \fmax{}, all      & 176.7 & \textbf{192.4} \\
\quad interpolation        & 198.7 & \textbf{209.3} \\
\quad extrapolation        & 138.6 & \textbf{163.0} \\
correctness                & 86.9\% & \textbf{95.1\%} \\
optimizer updates          & 284   & \textbf{200} \\
\midrule
\emph{conditional on correct}, interp. & 213.2 & 213.3 \\
\emph{conditional on correct}, extrap. & 182.1 & 181.6 \\
\bottomrule
\end{tabular}
\end{table}

The repaired reward is better on every axis, using $30\%$ fewer gradient
updates (Table~\ref{tab:v9}). The mechanism is visible in the last two rows:
frequency \emph{conditional on being correct} is unchanged to within
$0.1$\,MHz. Re-anchoring did not make designs faster. It stopped the policy
producing incorrect ones---and since the correctness regression was the
saturation event (Section~\ref{sec:trajectory}), removing the saturation removes
the regression.

Per family, the largest improvement is on \texttt{iir}
($129.2 \rightarrow 164.6$\,MHz, $+27\%$), the family whose predictions were
worst before repair. This is the prediction made in
Section~\ref{sec:divergence}, and it holds. The \texttt{med} family is unchanged
($20.0 \rightarrow 18.4$): a combinational sorting network's critical path is
not shortened by coding style, which is a limitation of the design corpus rather
than of the reward.

% ===========================================================================
\section{Evaluation Integrity}
\label{sec:integrity}
% ===========================================================================

Because our claims rest on an executable oracle, we audited the acceptance rule
rather than asserting it. All $1{,}529$ accepted candidates were re-scored under
three stricter rules: exact equality rather than the deployed threshold;
comparison of the post-reset window that the deployed rule skips; and a latency
offset inferred on one seed and frozen for the other. \emph{None} of the
$1{,}529$ was rejected by any of the three.

The deployed threshold is near-exact by construction: it compares
$k = n - \ell - w$ samples, so at the training length it admits zero mismatches
and at the evaluation length at most one in $1016$. We report the sample counts
rather than the threshold, since the latter reads as an approximate-similarity
score and is not one.

% ===========================================================================
\section{Limitations}
\label{sec:limitations}
% ===========================================================================

\textbf{One device, one tool, one surrogate architecture.} All measurements are
Vivado 2023.1 on a Zynq-7020 with a twelve-feature lexical predictor. Whether
the collapse threshold transfers to graph-based predictors or other flows is
untested.

\textbf{WNS-derived estimates, not timing closure.} As stated in
Section~\ref{sec:setup}. Comparisons are relative and use an identical flow
throughout; absolute frequencies should not be read as closed timing.

\textbf{Procedurally generated designs.} The thirty held-out designs come from
one parameterised catalogue, so family should be treated as the unit of
independence rather than design. With five families, statistical claims are
correspondingly weak, and we make none.

\textbf{Single seed.} The results reported here use one optimization seed per
configuration. Replication across seeds is in progress and will accompany the
full version.

\textbf{Saturation is partially, not wholly, repaired.} Six cells---the two
largest tap counts, across every policy including the supervised
one---still reach the clip bound after re-anchoring. The repair reduces the
failure; it does not eliminate it.

% ===========================================================================
\section{Related Work}
\label{sec:related}
% ===========================================================================

\textbf{Reinforcement learning for RTL quality.} ChipSeek-R1~\cite{chipseek2025}
trains with hierarchical correctness and PPA rewards; PPA-RTL~\cite{ppartl2025}
uses post-synthesis PPA directly as the reinforcement signal. Both establish
that physical feedback can drive RTL generation, and we do not claim that
contribution. Our subject is orthogonal and, we argue, prior to theirs: the
conditions under which such a reward remains \emph{valid} while it is being
optimized. Neither work measures its reward's error against ground truth on a
distribution created by its own optimization.

\textbf{Reward overoptimization.} Gao \emph{et al.}~\cite{gao2023scaling}
characterize the phenomenon for language-model reward models using a synthetic
gold reward model as ground truth. We measure it with a synthesis tool and a
physical device. The qualitative shape we observe---stability followed by
saturation---is consistent with theirs; the ability to attach real units to the
gap is what hardware adds.

\textbf{Evaluating RTL optimization.} RTL-OPT~\cite{lu2026rtlopt} provides a
benchmark for optimization-oriented RTL evaluation, and
FormalRTL~\cite{formalrtl2026} couples reference models to formal equivalence
checking. Our oracle is simulation-based and we make no formal claim; the audit
in Section~\ref{sec:integrity} is what supports the acceptance rule.

% ===========================================================================
\section{Conclusion}
% ===========================================================================

A learned physical reward can pass every offline check available to the
practitioner who trains it, and be wrong by $236$\,MHz precisely where
optimization drives the policy. The failure is abrupt, it is accompanied by a
functional-correctness regression that is the same event rather than a separate
one, and it is not visible in cross-validation---indeed, cross-validation ranks
the broken predictor above its repair.

The remedy is inexpensive: forty-three real measurements taken from the region
the optimizer actually reaches. What is not inexpensive is knowing that the
remedy is needed, and the practical recommendation we draw is to monitor the
reward's saturation during training, since it costs nothing and is the earliest
available signal that the objective has stopped meaning what it did.

\section*{Reproducibility}

Code, the frozen held-out design set, all generated candidates, and every
post-implementation measurement underlying the tables above are available at
\url{https://github.com/MoAmineHallam/FPGA}. Each table regenerates from the
committed artifacts with a single named script.

\bibliographystyle{IEEEtran}
\bibliography{refs}

\end{document}
